%%
%% Author: shoupu_wan
%% 2019-04-13
%%
\appendix{Hoare's Partitioning Functions}\label{apx:impl-varieties}

\lstinputlisting[label={lst:part_nested_loop_var1},language=C,
caption={Partition scheme for quicksort variation 1.
For elements equal to the pivot, we don't stop '<=' or '<' is not critical.
But '<=' brings more saving on unnecessary swaps.}]
{part_nested_loop_var1.c}

\lstinputlisting[label={lst:part_nested_loop_var2},language=C,
caption={Partition scheme for quicksort variation 2.
Found another variation of partition scheme based on part-nested-loop
Diff wise, it lies between part-nested-loop and part-nested-loop-var
- the order of the two while loops is identical to part-nested-loop
- the boundary checking is identical to part-nested-loop-var
- the return value is shifted toward the left by one unit
}]
{part_nested_loop_var2.c}

\lstinputlisting[label={lst:part_nested_loop_var3},language=C,
caption={Partition scheme for quicksort variation 3.
Variation of partition scheme.
The two differs in three ways
- order of the two while loops
- increment condition for equal element pointer
- return pointer choice between s and e
}]
{part_nested_loop_var3.c}

\lstinputlisting[label={lst:part_nested_loop_var4},language=C,
caption={Partition scheme for quicksort variation 4.
Find yet another variation of implementation where the second nested while-loop
is made post-incremental.
This will remove the difficulty encountered in loop-thinning.}]
{part_nested_loop_var4.c}

All these listings are to be understood with some pivot-selection mechanism to avoid pivot
skewness.

\appendix{Quicksort Partitioning function in Go}\label{apx:partition-go}
